% Generated draft: requires mathematical and source review.
\documentclass[11pt]{article}
\usepackage[margin=25mm]{geometry}
\usepackage{fontspec}
\setmainfont{Latin Modern Roman}
\usepackage{amsmath,amssymb,mathtools,unicode-math}
\setmathfont{Latin Modern Math}
\usepackage{enumitem,needspace,xurl,hyperref}
\setlength{\parindent}{0pt}
\setlength{\parskip}{5pt}
\setlength{\emergencystretch}{3em}
\newcommand{\N}{\mathbb{N}}
\newcommand{\Z}{\mathbb{Z}}
\newcommand{\Q}{\mathbb{Q}}
\newcommand{\R}{\mathbb{R}}
\newcommand{\C}{\mathbb{C}}
\newcommand{\F}{\mathbb{F}}
\newcommand{\A}{\mathbb{A}}
\newcommand{\PP}{\mathbb{P}}
\newcommand{\E}{\mathbb{E}}
\newcommand{\eps}{\varepsilon}
\newcommand{\dd}{\,\mathrm{d}}
\begin{document}
\section*{UnsolvedMath ID: 30002115}
\textbf{Source problem code:} catalogue entry
\subsection{Definitions and mathematical statement}
Let \(\Omega\subset\R^m\) be a bounded, sufficiently smooth domain. For a standard spectral interpretation of the supplied statement, let \(-\Delta_D\) denote the positive Dirichlet Laplacian on \(\Omega\), and write its eigenvalues, counted with multiplicity, as
\[
0<\lambda_1(\Omega)\leq\lambda_2(\Omega)\leq\cdots.
\]
Choose an orthonormal basis \((u_j)_{j\geq1}\) of \(L^2(\Omega)\) satisfying
\[
-\Delta u_j=\lambda_j u_j\quad\text{in }\Omega,
\qquad u_j|_{\partial\Omega}=0.
\]
Here \(\langle h,u_j\rangle=\int_\Omega h\,\overline{u_j}\,\dd x\). Interpret an admissible harmonic test function as \(h\in L^2(\Omega)\) with \(\Delta h=0\) in the distributional sense. No boundary condition on \(h\) is supplied. The normalization, integrability, and harmonicity conventions above make the expression meaningful, but are not explicitly specified in the source record.

For \(\lambda\notin\{\lambda_j:j\geq1\}\), define
\[
Q_\lambda(h)=\sum_{j=1}^{\infty}
\frac{\lambda_j}{\lambda_j-\lambda}
\bigl|\langle h,u_j\rangle\bigr|^2.
\]
Under these conventions this series is absolutely convergent. The well-defined, nonresonant portion of the supplied question is: for every such \(\Omega\) and every
\[
\lambda\in(\lambda_1(\Omega),\infty)
\setminus\{\lambda_j:j\geq1\},
\]
does there exist a harmonic \(h\in L^2(\Omega)\) such that \(Q_\lambda(h)<0\)?

\textbf{Unresolved qualifications.} The supplied statement actually quantifies over every \(\lambda>\lambda_1(\Omega)\), including Dirichlet eigenvalues. At those parameters its displayed fractions are undefined, and no limiting convention, spectral projection, or other interpretation is supplied. Thus the nonresonant formulation above covers only the meaningful portion of that quantifier and does not resolve the full stated question. The admissible dimensions and the meaning of sufficiently smooth are also unspecified. The supplied background reports a one-dimensional failure and refers to a higher-dimensional setting, without explicitly stating the intended dimension restriction.
\subsection{Short English statement}
For each admissible domain and each parameter above its first Dirichlet eigenvalue, the question asks whether a square-integrable harmonic function can make the displayed weighted sum of squared eigenfunction coefficients negative. The sum has a standard interpretation away from the Dirichlet spectrum. Its interpretation at eigenvalues, the intended dimensions, and the precise regularity assumptions remain unresolved. The supplied open label is dated curation, not an independently verified status.
\subsection{Sources}
\begin{itemize}
\item Supplied local UnsolvedMath registry record, permanent ID \(30002115\), source problem code \texttt{catalogue entry}, title \emph{Harmonic Test Inequalities Above the First Eigenvalue}; collection: Oberwolfach Reports Open Problems. The record lists proposed year 2012 and names no individual proposer.
\item Supplied report citation: \emph{Geometric Aspects of Spectral Theory (2013)}, Oberwolfach Reports, pp.\ 2013--2076, DOI \texttt{10.4171/owr/2012/33}. \url{https://doi.org/10.4171/owr/2012/33}. The supplied publication title/year and DOI are retained without reconciling their year references.
\item Supplied literature source, labelled OWR report: \url{https://ems.press/content/serial-article-files/46403}.
\item Literature checked at: \texttt{2026-08-21}. This is supplied dated curation, not independently verified here. Its assessment labels the problem open, describes harmonic and lambda-harmonic inequalities as conjectures, reports a one-dimensional failure, and says that no proof in the stated higher-dimensional setting was found. The record calls this a dated literature triage and requests independent verification of time-sensitive status and cited claims.
\item The supplied statement status is \texttt{corrected\_verified}. Its verification note claims that the official report was checked, the complete statement confirmed, and previously undelimited notation normalized. These are supplied provenance claims; no source was independently visited for this draft.
\end{itemize}
\end{document}
